\section*{第 \ref{ch:g11:scal} 章 --- 平面上的数量积}

\begin{solution}{exo:g11:scal:1}
$(2,5)\cdot(3,-1) = 6 - 5 = 1$.

$(1,-4)\cdot(8,2) = 8 - 8 = 0$（正交向量）。

$\vec u\cdot\vec v = 3 \times 4 \times \cos\frac\pi3 = 12 \times \frac12
= 6$.
\end{solution}

\begin{solution}{exo:g11:scal:2}
$\vect{AB}\cdot\vect{AC} = 6 \times 6 \times \cos 60^\circ = 18$.

向量 $\vect{AB}$ 与 $\vect{BC}$ 的夹角为
$180^\circ - 60^\circ = 120^\circ$（将二者起点置于 $B$ 处），于是
$\vect{AB}\cdot\vect{BC} = 36 \cos 120^\circ = -18$.
\end{solution}

\begin{solution}{exo:g11:scal:3}
$\vec u \cdot \vec v = 3(t - 8) + t = 4t - 24 = 0$ 得 $t = 6$。

$\vec u - \vec v = (t - 4,\ 0)$，于是
$\vec u \cdot (\vec u - \vec v) = t(t-4) + 0 = t^2 - 4t$，它在
$t = 0$ 或 $t = 4$ 时为零。
\end{solution}

\begin{solution}{exo:g11:scal:4}
$\vec u \cdot \vec v = 3 + 2 = 5$, $\norm{\vec u} = \sqrt5$,
$\norm{\vec v} = \sqrt{10}$，于是
\[
\cos\theta = \frac{5}{\sqrt5\,\sqrt{10}} = \frac{5}{5\sqrt2}
= \frac{\sqrt2}{2},
\]
因此恰有 $\theta = 45^\circ$。
\end{solution}

\begin{solution}{exo:g11:scal:5}
$a^2 = 16 + 36 - 2 \times 4 \times 6 \times \cos\frac{2\pi}{3}
= 52 - 48 \times \left(-\frac12\right) = 76$，于是 $a = \sqrt{76} =
2\sqrt{19} \approx 8.7$.

对于第二个三角形，由余弦定理解出角的余弦，得
\[
\cos\widehat A = \frac{b^2 + c^2 - a^2}{2bc}
= \frac{25 + 16 - 49}{2 \times 5 \times 4} = -\frac{8}{40} = -\frac15,
\]
负值，故 $\widehat A$ 为钝角。
\end{solution}

\begin{solution}{exo:g11:scal:6}
$\vect{AB}\,(4, -2)$ 与 $\vect{AC}\,(2, 4)$：
$\vect{AB}\cdot\vect{AC} = 8 - 8 = 0$，于是 $A$ 处的角为直角。
各边长：$AB = \sqrt{16 + 4} = \sqrt{20}$，$AC = \sqrt{4 + 16} =
\sqrt{20}$，又 $\vect{BC}\,(-2, 6)$ 得 $BC = \sqrt{40}$。验证：
$BC^2 = 40 = 20 + 20 = AB^2 + AC^2$，此即余弦定理
在 $\cos\widehat A = 0$ 时的情形（毕达哥拉斯定理）。
\end{solution}

\begin{solution}{exo:g11:scal:7}
圆心--半径式：$(x + 2)^2 + (y - 3)^2 = 25$。

对于第二个圆，配方：
\[
x^2 - 6x + y^2 + 2y = 6
\iff (x - 3)^2 - 9 + (y + 1)^2 - 1 = 6
\iff (x-3)^2 + (y+1)^2 = 16 :
\]
圆心 $(3, -1)$，半径 $4$。
\end{solution}

\begin{solution}{exo:g11:scal:8}
$M(x,y)$ 在圆上当且仅当 $\vect{MA}\cdot\vect{MB} = 0$，其中
$\vect{MA}\,(-3 - x,\ -y)$ 与 $\vect{MB}\,(-x,\ 2 - y)$：
\[
(-3 - x)(-x) + (-y)(2 - y) = x^2 + 3x + y^2 - 2y = 0 .
\]
对于原点：$0 + 0 + 0 - 0 = 0$，故 $O$ 在圆上 ---
从几何上看，$\vect{OA}\,(-3,0)$ 与 $\vect{OB}\,(0,2)$ 正交，
于是从 $O$ 看线段 $[AB]$ 的张角为直角。（配方后：
圆心 $\left(-\frac32, 1\right)$，半径 $\frac{\sqrt{13}}{2}$。）
\end{solution}

\begin{solution}{exo:g11:scal:9}
过 $P$ 的垂线的法向量取 $d$ 的方向向量，
即 $(-1, 2)$，故方程为 $-x + 2y + c = 0$；过
$P(1,3)$ 得 $-1 + 6 + c = 0$，$c = -5$，于是 $x - 2y + 5 = 0$。

与 $d$ 的交点：由 $2x + y = 7$ 得 $y = 7 - 2x$；代入，
$x - 2(7 - 2x) + 5 = 5x - 9 = 0$，于是 $x = \frac95$，$y = \frac{17}5$：
$H\left(\frac95, \frac{17}5\right)$。于是
$\vect{PH}\left(\frac45, \frac25\right)$，且
\[
PH = \sqrt{\frac{16}{25} + \frac{4}{25}} = \frac{\sqrt{20}}{5}
= \frac{2\sqrt5}{5} \approx 0.89 .
\]
\end{solution}

\begin{solution}{exo:g11:scal:10}
相加得
\[
\norm{\vec u + \vec v}^2 = \norm{\vec u}^2 + 2\vec u\cdot\vec v +
\norm{\vec v}^2
\quad\text{和}\quad
\norm{\vec u - \vec v}^2 = \norm{\vec u}^2 - 2\vec u\cdot\vec v +
\norm{\vec v}^2,
\]
交叉项相消，于是
$\norm{\vec u + \vec v}^2 + \norm{\vec u - \vec v}^2 = 2\norm{\vec u}^2 +
2\norm{\vec v}^2$。在以 $\vec u$ 与 $\vec v$ 为边的平行四边形中，
其对角线为 $\vec u + \vec v$ 与 $\vec u - \vec v$，两条对角线的
平方和等于四条边的
平方和。
\end{solution}

\begin{solution}{exo:g11:scal:11}
\emph{1.} 由于 $I$ 为中点，$\vect{IB} = -\vect{IA}$，且
$IA = 2$。于是
\[
\vect{MA}\cdot\vect{MB}
= (\vect{MI} + \vect{IA})\cdot(\vect{MI} - \vect{IA})
= \norm{\vect{MI}}^2 - \norm{\vect{IA}}^2
= MI^2 - 4 .
\]

\emph{2.} 条件 $\vect{MA}\cdot\vect{MB} = 12$ 化为
$MI^2 = 16$，即 $MI = 4$：所求集合为以 $I$ 为圆心的
半径为 $4$ 的圆。
\end{solution}

\begin{solution}{pb:g11:scal:1}
\textbf{1.} $\vec u \cdot \vec v = -3 + 8 = 5$；
$\norm{\vec u} = 5$，$\norm{\vec v} = \sqrt5$；
$\cos\theta = \frac{5}{5\sqrt5} = \frac{1}{\sqrt5}$：
$\theta \approx 63.4^\circ$。

\textbf{2.} $6 - 4k = 0$：$k = \frac32$。

\textbf{3.} $W = 50 \times 10 \times \cos 60^\circ = 250$
焦耳：只有沿运动方向的分量做功。

\textbf{4.} $c^2 = 25 + 49 - 2 \times 35 \times \frac12 = 39$：
$c = \sqrt{39}$。

\textbf{5.} $\cos\gamma = \frac{16 + 36 - 64}{2 \times 4
\times 6} = -\frac14$：结果为负，故该角为钝角：
$\gamma \approx 104.5^\circ$。

\textbf{6.} 由坐标计算：$\vec u \cdot \vec v = \cos a \cos b +
\sin a \sin b$。几何法：两个单位向量，
$\vec u \cdot \vec v = 1 \times 1 \times \cos(a - b)$。
同一数的两种计算结果相等：
$\cos(a - b) = \cos a \cos b + \sin a \sin b$。

\textbf{7.} $\cos(a + b) = \cos(a - (-b)) = \cos a \cos(-b) +
\sin a \sin(-b) = \cos a \cos b - \sin a \sin b$。

\textbf{8.} $\sin(a + b) = \cos\left(\frac\pi2 - a - b\right)
= \cos\left(\left(\frac\pi2 - a\right) - b\right)
= \cos\left(\frac\pi2 - a\right)\cos b +
\sin\left(\frac\pi2 - a\right)\sin b
= \sin a \cos b + \cos a \sin b$。

\textbf{9.} $\cos 15^\circ = \cos 45^\circ \cos 30^\circ +
\sin 45^\circ \sin 30^\circ = \frac{\sqrt2}{2} \cdot
\frac{\sqrt3}{2} + \frac{\sqrt2}{2} \cdot \frac12 =
\frac{\sqrt6 + \sqrt2}{4} \approx 0.966$。又
$\sin 75^\circ = \cos 15^\circ$：数值相同。

\textbf{10.} 在第 7 题和第 8 题中取 $b = a$：
$\cos 2a = \cos^2 a - \sin^2 a = 2\cos^2 a - 1 = 1 - 2\sin^2 a$
（利用 $\cos^2 + \sin^2 = 1$），且
$\sin 2a = 2 \sin a \cos a$。

\textbf{11.} $\sin^2 15^\circ = \frac{1 - \cos 30^\circ}{2}
= \frac{2 - \sqrt3}{4}$，故
$\sin 15^\circ = \frac{\sqrt{2 - \sqrt3}}{2}$。将
$\frac{\sqrt6 - \sqrt2}{4}$ 平方：
$\frac{6 - 2\sqrt{12} + 2}{16} = \frac{8 - 4\sqrt3}{16} =
\frac{2 - \sqrt3}{4}$：同一个数的两种写法。

\textbf{12.} $\vect{AP}$ 可分解为沿 $d$ 方向的一部分
（与 $\vec n$ 数量积为零）及垂直于 $d$ 的另一部分，其
长度即所求距离；与单位法向量
$\frac{\vec n}{\norm{\vec n}}$ 作数量积，消去第一部分，
并度量第二部分。设 $A$ 在 $d$ 上：
$\vect{AP} \cdot \vec n = a x_0 + b y_0 - (a x_A + b y_A) =
a x_0 + b y_0 + c$，由此得公式。例如：
$\frac{\abs{15 + 24 - 12}}{5} = \frac{27}{5} = 5.4$。

\textbf{13.} 半径 $= 5.4$：
$(x - 5)^2 + (y - 6)^2 = 29.16$。

\textbf{14.} 距离 $= \frac{\abs{-12}}{5} = 2.4$。垂足：过原点的
垂线为 $(3t, 4t)$；代入直线：
$9t + 16t = 12$，$t = \frac{12}{25}$：垂足为
$\left(\frac{36}{25}, \frac{48}{25}\right)$，距原点的距离为
$5t = \frac{60}{25} = 2.4$。公式与
直接计算结果一致。

\textbf{15.} 底 $AB = 6$，高 $=$ 点 $C$ 到
直线 $(AB)$（即横轴 $y = 0$）：$5$；面积为 $15$。直线 $(AC)$：
$5x - 2y = 0$，且
$\operatorname{dist}(B, (AC)) = \frac{\abs{30}}{\sqrt{29}}$；
于是 $\frac12 \times AC \times \frac{30}{\sqrt{29}} =
\frac12 \times \sqrt{29} \times \frac{30}{\sqrt{29}} = 15$：
面积相同，取哪条高均可。

\textbf{16.} $\operatorname{dist} = \frac{\abs{2 + 1 - 6}}
{\sqrt2} = \frac{3}{\sqrt2} \approx 2.12 > 2$：直线
与圆相离 --- 即无公共点。

\textbf{17.} 引入 $I$：
$\vect{MA} \cdot \vect{MB} = (\vect{MI} + \vect{IA}) \cdot
(\vect{MI} + \vect{IB}) = MI^2 + \vect{MI} \cdot (\vect{IA} +
\vect{IB}) + \vect{IA} \cdot \vect{IB}$。中间项消去
（$\vect{IA} + \vect{IB} = \vec 0$），且
$\vect{IA} \cdot \vect{IB} = -IA^2 = -\frac{AB^2}{4}$
（长度为 $\frac{AB}2$ 的相反向量）。

\textbf{18.} $\vect{MA} \cdot \vect{MB} = 0 \iff MI^2 =
\frac{AB^2}{4} \iff MI = \frac{AB}{2}$：圆心为
$I$，半径为 $\frac{AB}2$ 的圆 --- 即以
$[AB]$ 为直径的圆。

\textbf{19.} 以 $H$ 为起点展开每一对（$\vect{BC} =
\vect{HC} - \vect{HB}$，等）：
\[
\vect{HA} \cdot \vect{HC} - \vect{HA} \cdot \vect{HB}
+ \vect{HB} \cdot \vect{HA} - \vect{HB} \cdot \vect{HC}
+ \vect{HC} \cdot \vect{HB} - \vect{HC} \cdot \vect{HA} = 0 :
\]
两两相消 --- 此恒等式
对\emph{任意}四点均成立。对于我们的 $H$，前两个数量积
依题设均为零，故
$\vect{HC} \cdot \vect{AB} = 0$：$H$ 在
过 $C$ 的高上。三条高交于一点。

\textbf{20.} 加法公式来自对同一数量积的两种计算：
一种取\emph{坐标}形式，另一种取
\emph{模与夹角}形式；距离公式来自
\emph{投影}形式；垂心的存在性来自
\emph{双线性}（沙勒关系与展开）形式。同一量计算两次，
两式相等，便得一条定理 --- 此即数量积的
全部本领。
\end{solution}
